\section{Proofs}
\label{sec:proofs}



\subsection{Proof of Proposition \ref{thm:Mn-Leb}}

This is already in Wang and Ramdas, just included here for completeness.\fc{todo}


\subsection{Proof of Theorem \ref{thm:recovery}}



\begin{proof}
We recall that
$$
h_{n}(\theta,u)=e^{-n\theta^2/2 }A_n(u \theta ).
$$
Write $M_{n,m}(u)=\int h_n(\theta,u)\Pi(\dd\theta)$ and $M_{n,m}^\Leb(u)=\int h_n(\theta,u)\dd\theta$. 

\emph{Step 1 (ratio convergence).}  Fix $R>0$ and split
\[
  \frac{M_{n,m}(u)}{a_m}
  \;=\;\underbrace{\int_{|\theta|\leq R}}_{=:I_1(R,m)}
   +\underbrace{\int_{R<|\theta|\leq m/2}}_{=:I_2(R,m)}
   +\underbrace{\int_{|\theta|>m/2}}_{=:I_3(m)}\;
   h_{n}(\theta,u)\,\frac{\pi(\theta-m)}{a_m}\,\dd\theta.
\]
We have, for any $\theta$,
$$
\pi(\theta-m)\sim a_m \sim m^{-\beta}L_-(m)
$$
as $m\to\infty$. By the uniform convergence theorem for regularly varying functions
\citep[Theorem~1.2.1]{Bingham1987}, 
$$
\sup_{|\theta|\leq R}|\pi(\theta-m)/a_m-1|\to 0
$$
as $m\to\infty$, so 
$$
I_1(R,m)\to\int_{|\theta|\leq R}h_{n}(\theta,u)\,\dd\theta.
$$
For $I_2$, on $|\theta|\leq m/2$ we have $|\theta-m|\in[m/2,\,3m/2]$, so
regular variation yields a constant $C_\star<\infty$ (independent of $R$ and of
$m$ for all large $m$) such that $\pi(\theta-m)/a_m\leq C_\star$; hence
$$
I_2(R,m)\leq C_\star\int_{|\theta|>R}h_{n}(\theta,u)\,\dd\theta,
$$ which vanishes
as $R\to\infty$. For $I_3$, using the fact that $\pi$ is a proper prior
(i.e.\ $\int_\mathbb{R}\pi=1$),
\[
I_3(m)\;=\;\int_{|\theta|>m/2}h_{n}(\theta,u)\,\frac{\pi(\theta-m)}{a_m}\,\dd\theta
   \;\leq\; a_m^{-1}\sup_{|\theta|>m/2}h_{n}(\theta,u),
\]
the supremum is exponentially small in $m$ by \eqref{eq:hu-env}, while $a_m$ is
only polynomially small, so $I_3(m)\to 0$.

Combining these three bounds, for each fixed $R>0$,
\begin{align*}
   \int_{|\theta|\leq R}h_{n}(\theta,u)\,\dd\theta
   \;\leq\; \liminf_{m\to\infty}\frac{M_{n,m}(u)}{a_m}
   &\;\leq\; \limsup_{m\to\infty}\frac{M_{n,m}(u)}{a_m}\\
   & \;\leq\; \int_{|\theta|\leq R}h_{n}(\theta,u)\,\dd\theta
   \;+\; C_\star\!\int_{|\theta|>R}h_{n}(\theta,u)\,\dd\theta.
\end{align*}
Sending $R\to\infty$ and using $\int_\mathbb{R} h_{n}(\theta,u)\,\dd\theta
= M_n^{\Leb}(u)$ gives \eqref{eq:Mnm-asymp}.

\emph{Step 2 (calibration convergence).}
Since \(|U_2|<\sqrt 2\) almost surely under \(\Prob_0\), Step~1 applied with \(n=2\) gives
\(M_{2,m}(U_2)/a_m\to M_2^{\Leb}(U_2)\) almost surely. Since \(0\leq\min(M_{2,m}(U_2)/(a_m c),1)\leq 1\), bounded convergence yields
\[
  g_m(a_m c)
  \;=\;\E_0\!\Bigl[\min\!\Bigl(\tfrac{M_{2,m}(U_2)}{a_mc},1\Bigr)\Bigr]
  \;\longrightarrow\;
  \E_0\!\Bigl[\min\!\Bigl(\tfrac{M_2^{\Leb}(U_2)}{c},1\Bigr)\Bigr]
  \;=\;g_{\Leb}(c),
\]
proving \eqref{eq:gm-asymp}.

\emph{Step 3 (composition).}
Set \(c_m:=M_{n,m}(u_n)/a_m\) and \(c:=M_n^{\Leb}(u_n)\); by Step~1 \(c_m\to c\). Each \(g_m\) is decreasing, so for any \(\varepsilon\in(0,c)\) and \(m\) large, \(g_m(a_m(c+\varepsilon))\leq g_m(a_m c_m)\leq g_m(a_m(c-\varepsilon))\). Taking \(\liminf\) and \(\limsup\) using Step~2 and letting \(\varepsilon\downarrow 0\) yields \eqref{eq:p-asymp}.
\end{proof}




\section{Secondary lemmas}
\label{sec:secondarylemmas}

\begin{lemma}\label{lemma:hthetabound}
For $n\geq 2$, $|u|<\sqrt{n}$, $\theta\in\bbR$, let 
  \begin{equation}
\label{eq:hthetadef2}
  h_{n}(\theta,u)
  =
  e^{-n\theta^2/2}
  A_n\bigl(\theta u\bigr),
\end{equation}
where $A_n$ is the MGF of $\chi_n$.  Then, for any $u^2<\eta<n$, 
\begin{equation}
h_{n}(\theta,u)\leq \exp\left(-\frac{(n-\eta)\theta^2}{2}\right)\left(1-\frac{u^2}{\eta}\right)^{-n/2}.
\label{eq:hu-env}
\end{equation}
Thus, there are constants $C_{n,u}$,$\kappa_{n,u}>0$ such that 
$$
0\leq h_{n}(\theta,u)\leq C_{n,u}e^{-\kappa_{n,u}\theta^2}.
$$
\end{lemma}

\begin{proof}
\fc{TODO}
\end{proof}


\begin{lemma}\label{sec:asymp_An}
For $R_n\sim\chi_n$, let
\[
A_n(z)=\mathbb E[e^{z R_n}]=\frac{1}{\Gamma(n/2)}\int_0^\infty y^{n/2-1}\exp(-y+z\sqrt{2y})dy.
\]
Then
$$A_n(z) \sim \frac{2^{1-n/2}}{\Gamma(n/2)} \sqrt{2\pi} z^{n-1} e^{z^2/2}$$ as $z \longrightarrow \infty$.
\end{lemma}











\begin{proof}
Substituting \(r=\sqrt{2y}\), so that \(y=r^2/2\) and
\(dy=r\,dr\), and completing the square, we obtain
\[
A_n(z)
=
\frac{2^{1-n/2}}{\Gamma(n/2)}
e^{z^2/2}
\int_0^\infty
r^{n-1}e^{-(r-z)^2/2}\,dr.
\]
Setting \(u=r-z\) gives
\[
\int_0^\infty
r^{n-1}e^{-(r-z)^2/2}\,dr
=
z^{n-1}
\int_{-z}^\infty
\left(1+\frac{u}{z}\right)^{n-1}
e^{-u^2/2}\,du.
\]
For each fixed \(u\in\mathbb{R}\),
\[
\left(1+\frac{u}{z}\right)^{n-1}
\mathbf{1}_{\{u>-z\}}
\longrightarrow 1
\qquad \text{as } z\to\infty.
\]
Moreover, for \(z\geq 1\), the integrand is bounded by
\[
(1+|u|)^{n-1}e^{-u^2/2},
\]
which is integrable on \(\mathbb{R}\). Therefore, by the dominated
convergence theorem,
\[
\int_{-z}^\infty
\left(1+\frac{u}{z}\right)^{n-1}
e^{-u^2/2}\,du
\longrightarrow
\int_{-\infty}^\infty e^{-u^2/2}\,du
=
\sqrt{2\pi}.
\]
Consequently,
\[
A_n(z)
\sim
\frac{2^{1-n/2}}{\Gamma(n/2)}
\sqrt{2\pi}\,
z^{n-1}e^{z^2/2},
\qquad z\to\infty.
\]
\end{proof}